Holographic near-eye displays are an emerging class of 3D display technology that enable realistic, immersive visual experiences through physically accurate wavefront reconstruction \cite{maimone2017holographic, chao2024large, gopakumar2024full, kim:HolographicGlasses, jang2024waveguide}. However, despite their immense promise, their practical deployment remains fundamentally constrained by the extreme computational cost of computer-generated holography (CGH). Driving such displays requires synthesizing interference patterns from complex 3D inputs in real time, yet this mapping from scene representation to SLM wavefront modulation is computationally demanding.

Existing CGH algorithms work with a wide range of 3D representations. However, each representation introduces significant computational trade-offs that limit real-time CGH:
RGBD images \cite{choi2021neural, peng2020neural} suffer from occlusion artifacts, mesh- \cite{matsushima2005computer, matsushima2009extremely}, point-cloud- \cite{h2009computer, maimone2017holographic}, or Gaussian-based \cite{choi2025gaussian, chao2025rpws} pipelines must process millions of primitives, and densely-sampled light fields \cite{chao2024large, schiffers2023stochastic} require content that is rarely produced in typical capture settings. As a result, generating high-quality, photorealistic holograms at interactive rates from compact or casually captured inputs remains an open challenge. 

Multiplane images (MPIs) \cite{zhou2018mpi, pu2023sinmpi, khan2023tiledmpi, han2022adaptivempi, tucker2020singlempi, flynn2019deepviewmpi, zhao2022generativempi, luvizon2021adaptivempi, srinivasan2019pushingmpi} are a layered 3D representation that consists of a set of RGBA images, each defined at a corresponding depth which can be predetermined or adaptively optimized from a set of casually captured images of a front-facing 3D scene. Novel views are synthesized through warping the RGBA images from a reference viewpoint to the novel viewpoint and subsequent alpha compositing. Although a single set of MPIs is unable to achieve 360-degree, free-view 3D reconstruction due to its fronto-parallel arrangement, it provides exceptional reconstruction of novel views particularly for \textit{local, front-facing light fields} of a 3D scene. This property is especially well-suited for CGH since only front-facing views need to be reconstructed well in the small eyebox of a near-eye display. Most importantly, MPIs usually only consist of a handful of RGBA images (<100), allowing for much faster processing in the wave optics domain compared to Gaussians and meshes where millions of primitive wavefronts have to be propagated and composited \cite{choi2025gaussian, matsushima2009extremely, matsushima2018full}.

In this paper, we devise a CGH pipeline that uses MPIs as an intermediate 3D representation for efficient computer-generated holography. We show that MPI-based CGH achieves high-quality focal stack and light field reconstruction quality and outperforms other conventional layer-based CGH methods such as CGH using RGBD images. We further demonstrate that MPI-based CGH can achieve up to $250,000\times$ speedup compared to state-of-the-art primitive-based methods such as polygon- or Gaussians-based CGH \cite{matsushima2005computer, matsushima2009extremely, matsushima2014silhouette, choi2025gaussian}, paving the way for practical and real-time CGH.

Specifically, our contributions include:
\begin{itemize}
\item A novel CGH algorithm that composites multiplane images into complex holograms with optional time multiplexing.
\item A detailed evaluation of MPI-based holograms against existing CGH baselines, analyzing how the number of depth layers impacts image quality and runtime.
\item Experimental validation demonstrating accurate 3D focal stack reconstruction and \textit{4D light field parallax} in MPI-based holograms.
\end{itemize}

% mathematical analysis of the depth binning

Source code and example datasets will be made public.
